\documentclass[11pt,a4paper]{article}
\usepackage[margin=1in]{geometry}
\usepackage{hyperref}
\usepackage{enumitem}
\usepackage{graphicx}
\usepackage{xcolor}

% -----------------------------------------------------
% bvraghav-tiet-ucs503p-article.sty
% -----------------------------------------------------

% Variable: Lab Instructor
\makeatletter \newcommand{\BvrSetLabInstructor}[1]{%
  \gdef\Bvr@LabInstructor{#1}}
% default
\newcommand{\Bvr@LabInstructor}{Hardik Gujral}
% public accessor
\newcommand{\BvrLabInstructorValue}{\Bvr@LabInstructor}
\makeatother

% Add subtitle
\usepackage{titling}
% -- Set title bold
\pretitle{\begin{center}\bfseries\fontsize{18bp}{18bp}\selectfont}
\makeatletter
\newcommand{\BvrSubtitle}[1]{%
  \posttitle{%
    \par\end{center}
    \begin{center}\large#1\end{center}
    \vskip 0.5em}%
}
\makeatother

% Hints in the section title(s)
\newcommand{\BvrHint}[1]{%
  {\normalfont\normalsize\itshape\hfill #1}}

% -----------------------------------------------------

\title{LocalHands: Discover and Directly Connect with Local Service Providers}

\BvrSetLabInstructor{Hardik Gujral}

\BvrSubtitle{%
  Project Proposal for UCS503P  \\
  Submitted to: \BvrLabInstructorValue \\ \bfseries
  Thapar Institute of Engineering and Technology% 
}

\author{
  Pratham Goyal%
  {Roll No: \texttt{1024170233}, %
    Email: \texttt{pgoyal4\_be24@thapar.edu}}
  \and 
   Y Vijay Kiran%
  {Roll No: \texttt{1024170415}, %
    Email: \texttt{ykiran\_be24@thapar.edu}}
  \and 
  Keshav Goyal%
  {Roll No: \texttt{1024170318}, %
    Email: \texttt{kgoyal1\_be24@thapar.edu}}
}

\date{}

\begin{document}
\maketitle

\section{Higher-order goal%
  \BvrHint{\ldots secondary framing}}
This project aims to improve access to trusted local
labour services (electricians, plumbers, carpenters,
sweepers, painters, and similar informal-sector
workers) by making it easier for households to discover
and directly contact a nearby worker. While the
underlying theme is broader -- digitizing trust and
discoverability in the informal service economy -- the
immediate engineering objective is to translate this
into a measurable, deployable software solution: a
searchable directory of verified worker profiles that
customers can browse and call directly.

\section{Time-to-value %
\BvrHint{\ldots primary heuristic}}
\begin{itemize}[leftmargin=0.7cm]
    \item We will deliver meaningful increments quickly by prototyping the core search-to-call workflow first, before adding any secondary features.
    \item We will prioritize fast feedback over perfection by using weekly sprint iterations and validating the workflow with a small set of test worker profiles early.
    \item Success is defined by what can be delivered and measured in the first iteration -- a working search $\rightarrow$ profile $\rightarrow$ call flow -- then improved based on feedback.
\end{itemize}

\section{Problem Statement}
Households needing local labour (electricians,
carpenters, plumbers, sweepers, etc.) typically rely on
word-of-mouth or a single known contact, with no
organized way to discover or compare nearby workers.
Common pain points include:
\begin{itemize}[leftmargin=0.7cm]
    \item \textbf{No discoverability:} there is no centralized way to find multiple nearby workers for a given service category.
    \item \textbf{Low transparency:} customers cannot easily compare experience, certifications, or approximate charges before making contact.
    \item \textbf{Trust deficit:} customers have no way to verify whether a worker is genuine or capable before calling.
    \item \textbf{Inefficient search:} without a directory, customers spend time calling around or asking neighbours to find someone quickly.
\end{itemize}

\textbf{Why this matters now (contextual relevance):} With increasing smartphone penetration even in tier-2/3 cities and neighbourhoods, a simple directory-and-call model requires no behavioural change from either customers or workers -- the app only replaces the discovery step, not the actual transaction, which keeps adoption friction low.

\section{Proposed Solution %
\BvrHint{\ldots Software Proposition}}
\subsection{Overview}
Build a \textbf{web-based} ``Discover-and-Call'' directory platform with the following components:
\begin{itemize}[leftmargin=0.7cm]
    \item \textbf{Worker registration portal:} workers create a profile with name, phone number, service category, years of experience, optional certifications, approximate charges, and service area/location.
    \item \textbf{Admin verification:} an administrator reviews and approves a worker profile before it becomes publicly visible, as a basic trust/anti-fake-profile measure.
    \item \textbf{Customer search:} customers search by service category and location (pincode/area) and view a ranked list of nearby matching worker profiles.
    \item \textbf{Profile detail and direct contact:} customers view full profile details and tap a ``Call Now'' action to contact the worker directly by phone -- no in-app booking or payment is involved.
    \item \textbf{Post-interaction rating:} customers may optionally leave a rating/review after contacting a worker, to build trust signals over time.
\end{itemize}

\subsection{Core workflow}
\begin{enumerate}[leftmargin=0.7cm]
    \item Worker registers with profile details and is marked pending until admin approval.
    \item Admin reviews and approves (or rejects) the profile.
    \item Customer searches by service category and location; system returns the nearest 6--7 approved, matching worker profiles.
    \item Customer opens a profile, reviews experience/certification/charges, and taps ``Call Now'' to contact the worker directly.
    \item (Optional) Customer leaves a rating after the interaction.
\end{enumerate}

\subsection{Operational constraints for an educational, tier-2/3 setting}
\begin{itemize}[leftmargin=0.7cm]
    \item Keep the solution usable on low-to-mid range devices via a responsive, mobile-first web design, since both workers and customers are likely to access it primarily on phones.
    \item Avoid dependence on advanced research algorithms (no recommendation engine or ML matching required); focus on correct search, filtering, and verification workflow.
    \item Design for deployability using common web hosting and managed databases.
    \item Minimize personal data collected (only what is necessary for discovery and contact) to reduce privacy risk.
\end{itemize}

\section{Solution Approach %
\BvrHint{\ldots Engineering focus}}
This is an engineering course project, so we focus on
correct workflow, search, and verification rather than
novel algorithm research.

\subsection{Search and ranking %
\BvrHint{\ldots non-research}}
Matching and ranking will be heuristic-based, not ML-driven:
\begin{itemize}[leftmargin=0.7cm]
    \item Filter worker profiles by exact/nearby service category and matching location (pincode/area text match for the MVP).
    \item Rank results by proximity (distance if geolocation is available, otherwise same-area priority) and then by average rating.
    \item Limit results to the top 6--7 nearest matches per search to keep the customer's decision simple.
\end{itemize}

\subsection{Web architecture %
\BvrHint{\ldots web-based solution}}
A typical 3-tier web architecture:
\begin{itemize}[leftmargin=0.7cm]
    \item \textbf{Frontend:} responsive UI (React) for customer search/profile browsing and worker registration/dashboard.
    \item \textbf{Backend API:} authentication (worker/admin), profile CRUD, search/filter endpoints, rating endpoints (Node.js/Express).
    \item \textbf{Database:} store users, worker profiles, certifications, service categories, and reviews (MongoDB).
\end{itemize}

\subsection{CI/CD and fast delivery enabler}
CI/CD will be used to:
\begin{itemize}[leftmargin=0.7cm]
    \item Automatically build and run tests on every change.
    \item Produce repeatable deployment artifacts to staging (e.g., Render/Railway for backend, Vercel/Netlify for frontend).
    \item Enable safe and frequent releases of incremental features across weekly sprints.
\end{itemize}

\section{Evaluation Criterion %
\BvrHint{\ldots measurable and attributable}}
We define success using measurable metrics tied to the core search-to-call workflow.

\subsection{Primary evaluation metric}
\textbf{Time-to-Contact (TTC):} time between a customer submitting a search query and the customer tapping ``Call Now'' on a profile.
\begin{itemize}[leftmargin=0.7cm]
    \item Target: median TTC $\le$ 60 seconds during pilot testing.
    \item Attribution: measured from application logs/timestamps (search event vs. call-action event).
\end{itemize}

\subsection{Secondary evaluation metrics}
\begin{itemize}[leftmargin=0.7cm]
    \item \textbf{Search-to-call conversion rate:} percentage of searches that result in at least one ``Call Now'' tap.
    \item \textbf{Profile approval turnaround:} time between worker registration and admin approval.
    \item \textbf{Profile completeness rate:} percentage of approved worker profiles with certification and pricing fields filled in.
    \item \textbf{Reliability:} availability of the search and profile pages during testing (e.g., $\ge$ 99\% uptime in pilot).
\end{itemize}

\subsection{Pilot validation plan %
\BvrHint{\ldots high-level}}
\begin{itemize}[leftmargin=0.7cm]
    \item Onboard 10--15 sample worker profiles across 4--5 service categories (electrician, plumber, carpenter, sweeper, painter).
    \item Have 10--15 test users perform searches and report on ease of finding and deciding to call a worker, via a short feedback form.
\end{itemize}

\section{Scalability %
\BvrHint{\ldots with foresight}}
The proposition should scale in both theoretical and practical senses.

\begin{enumerate}[leftmargin=0.7cm]
    \item \textbf{Scalable (theoretically):} The system should support growth in the number of worker profiles, service categories, and concurrent searches by:
    \begin{itemize}[leftmargin=0.7cm]
        \item decoupling frontend and backend,
        \item indexing search fields (service category, location/pincode) for fast lookups,
        \item using stateless API design so multiple backend instances can be run behind a load balancer if needed.
    \end{itemize}
    
    \item \textbf{Operationally deployable (practically):} The system should run on common web hosting:
    \begin{itemize}[leftmargin=0.7cm]
        \item managed NoSQL database (MongoDB Atlas),
        \item platform-hosted deployment (Render/Vercel-class services),
        \item CI/CD pipeline for staging/production.
    \end{itemize}
\end{enumerate}

\section{Engine availability heuristic}
A mature set of ``engines'' (tooling/services) is available to implement this as a web-based project:
\begin{itemize}[leftmargin=0.7cm]
    \item Web framework for REST APIs (Express.js) and a reactive frontend (React.js).
    \item Managed document database (MongoDB Atlas) with geospatial/text indexing support.
    \item Authentication approach (JWT-based, token authentication for worker/admin roles).
    \item Object storage for profile photos and certification files (Cloudinary).
    \item Test framework for unit/integration tests (Jest).
    \item CI runner and staging deployment target (GitHub Actions, Render/Vercel).
\end{itemize}
We will use these mature components to focus on workflow design, search correctness, and measurement rather than inventing new infrastructure.

\section{Project scope and deliverables %
\BvrHint{\ldots iteration-friendly}}
\subsection{Initial deliverable %
\BvrHint{\ldots first iteration}}
\begin{itemize}[leftmargin=0.7cm]
    \item Worker registration form + profile schema (name, phone, category, experience, charges, location)
    \item Admin approval workflow (approve/reject pending profiles)
    \item Customer search by service category + location, returning nearest matches
    \item Worker profile detail page with ``Call Now'' action
    \item Basic logging and timestamps for Time-to-Contact measurement
    \item CI: build + backend unit tests + linting
    \item CD to staging with a single pipeline job
\end{itemize}

\subsection{Subsequent deliverables}
\begin{itemize}[leftmargin=0.7cm]
    \item Certification upload and display (badge on profile)
    \item Rating/review system for workers
    \item Geolocation-based distance search (upgrade from pincode/area text match)
    \item Worker dashboard: edit profile, view search/call impressions
    \item ``Chat on WhatsApp'' as an additional contact option alongside Call Now
\end{itemize}

\section{Risks and mitigations}
\begin{itemize}[leftmargin=0.7cm]
    \item \textbf{Fake or low-quality profiles:} mitigated by mandatory admin approval before a profile is publicly visible.
    \item \textbf{Phone number misuse/privacy:} minimize collected personal data; expose the phone number only on the profile detail page, not in search result previews.
    \item \textbf{Low digital literacy among workers:} keep the registration form short with minimal required fields; consider phone-assisted onboarding as a stretch goal.
    \item \textbf{No feedback loop after a call (since booking is outside the app):} address via an optional, lightweight post-interaction rating prompt rather than mandatory booking confirmation.
    \item \textbf{Evaluation measurement ambiguity:} instrument search and call-action timestamps and define clear event logging before development begins.
\end{itemize}

\section{Summary}
This project proposes a deployable web platform,
LocalHands, that reduces the friction of discovering
trustworthy local labour by replacing word-of-mouth
search with a searchable, verified worker directory. It
meets the course criteria by emphasizing time-to-value
through rapid iterations, implementing CI/CD to support
frequent safe releases, and using measurable evaluation
criteria (especially Time-to-Contact). By intentionally
scoping out in-app booking and payments, the project
remains focused on engineering a correct, scalable
search-and-discovery workflow rather than research-driven
novelty.

\end{document}